% Part: first-order-logic
% Chapter: introduction
% Section: syntax

\documentclass[../../../include/open-logic-section]{subfiles}

\begin{document}

\olfileid{fol}{int}{syn}

\olsection{Syntaxis}

We moeten eerst precies vastleggen welke reeksen symbolen als
!!{sentence}s van de eerste-ordelogica gelden. Dat doen we later;
voorlopig werken we met voorbeelden. De basiselementen---de
woordenschat---van de eerste-ordelogica vallen in twee delen uiteen.
Het eerste deel bestaat uit de symbolen waarmee we bepaalde dingen
aanduiden of er iets bepaalds over zeggen. We duiden dingen aan met
!!{constant}s en zeggen iets over de aangeduide dingen met
!!{predicate}s. Zo kunnen we $\Obj a$ als !!a{constant} gebruiken om
één ding aan te duiden en daar vervolgens met de
!!{sentence}~$\Atom{\Obj P}{\Obj a}$ iets over zeggen. Als men een
betekenis voor ``$\Obj a$'' en ``$\Obj P$'' voor ogen heeft, kan men
$\Atom{\Obj P}{\Obj a}$ als een Nederlandse zin lezen (en dat heeft
men waarschijnlijk gedaan toen men voor het eerst formele logica
leerde). Zodra men zulke eenvoudige !!{sentence}s van de
eerste-ordelogica heeft, kan men met het tweede deel van de
woordenschat---de logische symbolen (connectieven en kwantoren)---
complexere zinnen vormen. Zo kunnen we bijvoorbeeld uitdrukkingen
vormen als $(\Atom{\Obj P}{\Obj a} \land \Atom{\Obj Q}{\Obj b})$
of~$\lexists[\Obj x][\Atom{\Obj P}{\Obj x}]$.

Om de precieze definities van de semantiek en de regels van onze
systemen voor !!{derivation} te kunnen geven die voor rigoureus
metalogisch onderzoek nodig zijn, moeten we allereerst precies
definiëren wat als !!a{sentence} van de eerste-ordelogica geldt. Het
grondidee is eenvoudig: sommige eenvoudige !!{sentence}s, zoals
$\Atom{\Obj P}{\Obj a}$, kunnen we alleen uit !!{predicate}s en
!!{constant}s vormen. Daaruit vormen we met de connectieven en
kwantoren complexere zinnen. Maar volgens welke regels mogen we
precies complexere !!{sentence}s vormen? Die regels moeten worden
vastgelegd; anders is ``!!{sentence} van de eerste-ordelogica'' niet
nauwkeurig genoeg gedefinieerd. Daarbij doen zich enkele kwesties
voor. Ten eerste moeten precies de juiste reeksen symbolen als
!!{sentence}s gelden. Ten tweede moet de definitie zo zijn opgezet dat
we wiskundige bewijzen over \emph{alle} !!{sentence}s kunnen geven.
Ten slotte moeten we ook enkele elementaire bewerkingen op
!!{sentence}s precies definiëren, zoals ``vervang elke $\Obj x$ in~$!A$
door~$\Obj b$.'' Door de kwantoren en !!{variable}s van de
eerste-ordelogica is niet zonder meer duidelijk hoe dat moet gebeuren.
Moet bijvoorbeeld $\lexists[\Obj x][\Atom{\Obj P}{\Obj a}]$ als
!!a{sentence} gelden? En
$\lexists[\Obj x][\lexists[\Obj x][\Atom{\Obj P}{\Obj x}]]$? Wat moet
het resultaat zijn van ``vervang $\Obj x$ door~$\Obj b$ in
$(\Atom{\Obj P}{\Obj x} \land \lexists[\Obj x][\Atom{\Obj P}{\Obj
x}])$''?

\end{document}
